\section{Experimental Results}
\label{sec:results}
Of the Study~I hypotheses, H1$'$ (dose-response) passes; H2$'$ passes for bone-region PSNR against the reference model and against the
no-twin control and fails for SSIM; H3$'$ passes for both contrasts on VinDr-RibCXR; H4$'$a is mixed (ours above the reference model on
JSRT and below it on Node21, above the no-twin control on both); H4$'$b shows no downstream benefit of the soft-tissue image used alone but
gains when the suppressed images are added as extra channels; and H5$'$ holds on synthetic data with exact truth and fails under the shift
to real radiographs. The v1 hypothesis H1 passes, and the exploratory site calibration passes H16--H18. In Study~II the primary hypothesis
H6 fails on all three data sets and H7 finds no difference from the capacity control, whereas H8 finds a gain for focal and a loss for
diffuse findings on NIH and H10 passes; the post hoc recalibration of the bone intervals (H9) brings their coverage close to nominal.
Table~\ref{tab:hypotheses} lists every hypothesis with its main estimate and outcome; the subsections below give the details.

\begin{table*}[!t]
\centering
\caption{Hypotheses, their main estimate with 95\,\% interval (or one-sided bound) and outcome. Study~I: final suppressors (H1$'$--H5$'$),
v1 models (H1), site calibration (H16--H18, exploratory) and external validation (E1--E3); Study~II: fusion study (H6--H10; H9 concerns
the suppressor's intervals and is post hoc). $\Delta$: ours (or the first-named model) minus the comparator.}
\label{tab:hypotheses}
\footnotesize
\renewcommand{\arraystretch}{1.08}
\setlength{\tabcolsep}{3pt}
\begin{tabularx}{\textwidth}{@{}>{\raggedright\arraybackslash}p{4.3cm}>{\raggedright\arraybackslash}X%
>{\raggedright\arraybackslash}p{4.6cm}>{\raggedright\arraybackslash}p{2.3cm}@{}}
\toprule
Hypothesis & Data and test & Estimate & Outcome \\
\midrule
H1$'$ absorption falls with the twin weight & VinDr-RibCXR; slope of $A$ on $\log_2(1+w)$, image-cluster bootstrap & $-0.0233$
 $[-0.0249,-0.0217]$ & passed \\
H2$'$ bone-region PSNR$_{\mathrm{cal}}$, ours vs.\ reference & 51 held-out SZCH pairs; non-inferiority, margin 0.5\,dB & $-0.06$\,dB
 (lower bound $-0.15$\,dB) & passed \\
H2$'$ bone-region PSNR$_{\mathrm{cal}}$, ours vs.\ no-twin & same; margin 0.25\,dB & $+0.01$\,dB (lower bound $-0.02$\,dB) & passed \\
H2$'$ SSIM$_{\mathrm{cal}}$, ours vs.\ reference & same; margin 0.01 & $-0.021$ (lower bound $-0.023$) & failed \\
H3$'$ $A$, ours vs.\ reference / no-twin & VinDr-RibCXR, 1,470 sites; paired, Holm & $-0.098$ $[-0.108,-0.089]$ / $-0.127$
 $[-0.134,-0.119]$ & passed \\
H4$'$a $\Delta$CNR of real nodules, ours vs.\ reference & JSRT (154) / Node21 (1,474) & $+0.023$ $[0.008,0.041]$ / $-0.025$
 $[-0.033,-0.018]$ & mixed \\
H4$'$b downstream, our soft image alone & NIH macro-AUROC & $-0.001$ $[-0.004,0.003]$ & no gain \\
H4$'$b downstream, extra channels & NIH macro-AUROC / Node21 CPM & $+0.006$ $[0.003,0.008]$ / $+0.028$ $[0.009,0.047]$ & gain \\
H5$'$ conformal coverage (nominal 0.90) & Qure.ai test split / 51 held-out SZCH pairs & 0.900 $[0.897,0.903]$ / 0.139
 $[0.134,0.145]$ & held / failed under shift \\
H1 (v1) CLIA vs.\ recipe & synthetic LIDC-IDRI test CTs, exact truth & $-0.022$ $[-0.032,-0.013]$ & passed \\
H16 / H16b site-calibration gain of ours & 51 held-out SZCH pairs; one-sided bound & $+2.12$\,dB (bound 1.99) / $+2.74$\,dB (bound 2.61)
 & passed \\
H17 calibrated ours vs.\ calibrated reference & same; margin 0.5\,dB & $-0.35$\,dB (lower bound $-0.44$\,dB) & passed \\
H18 $A$, calibrated ours vs.\ calibrated reference & VinDr-RibCXR, 1,470 sites & $-0.096$ $[-0.105,-0.087]$ & passed \\
E1 pneumonia macro-AUROC, ours vs.\ original & BDCXR-3257, five folds; margin 0.01 & $+0.002$ $[-0.005,0.010]$ & non-inferior \\
E2 $A$, ours vs.\ reference & BDCXR-3257, 300 images & $-0.134$ $[-0.144,-0.124]$ & lower \\
E3 ensemble spread, children vs.\ adults & BDCXR-3257 vs.\ SZCH held-out & 2.7-fold (0.0076 vs.\ 0.0028) & descriptive \\
\midrule
H6 (primary) macro-AUROC, hybrid vs.\ original & NIH / VinBigData / BDCXR-3257; Holm & $+0.0001$ $[-0.0022,0.0024]$ / $+0.0024$
 $[-0.0004,0.0053]$ / $-0.0042$ $[-0.0108,0.0021]$ & failed \\
H7 hybrid vs.\ control / vs.\ early & NIH; Holm over six tests & $+0.0007$ $[-0.0018,0.0032]$ / $-0.0031$ $[-0.0055,-0.0007]$ & no
 difference \\
H8 focal findings, soft vs.\ original & NIH; Holm & $+0.0089$ $[0.0040,0.0140]$, Holm $p=0.0024$ & gain \\
H8 diffuse findings, hybrid vs.\ original & NIH; Holm & $-0.0052$ $[-0.0077,-0.0027]$, Holm $p=0.0024$ & loss \\
H9 recalibrated bone intervals (post hoc) & 24 development / 51 held-out SZCH pairs & coverage 0.886 (global), 0.895 (Mondrian) & near
 nominal \\
H10 pointing game, soft minus radiograph stream, focal vs.\ other & NIH, 984 expert boxes & $+0.145$ $[0.092,0.199]$, $p=0.0002$ &
 passed \\
\bottomrule
\end{tabularx}
\end{table*}

% ---------------------------------------------------------------------------------------------------------------------
\subsection{Fidelity, Dose-Response and Ablation (Study~I: H1\texorpdfstring{$'$}{'}, H2\texorpdfstring{$'$}{'})}
\label{sec:res-fidelity}

\subsubsection{Fidelity on held-out paired radiographs (H2$'$)}
On the 51 held-out SZCH-X-Rays pairs (Table~\ref{tab:fidelity}) the bone-region PSNR$_{\mathrm{cal}}$ of ours is 21.46\,dB and that of the
reference model 21.51\,dB. The paired difference is $-0.06$\,dB with a one-sided 95\,\% lower bound of $-0.15$\,dB, inside the
non-inferiority margin of $-0.5$\,dB. Against the no-twin control the difference is $+0.01$\,dB (lower bound $-0.02$\,dB, margin
$-0.25$\,dB): the twin regulariser costs no measurable bone-region fidelity. SSIM$_{\mathrm{cal}}$ fails its margin: ours is lower by 0.021
(lower bound $-0.023$, margin $-0.01$). The gap sits in fine texture. With a band limit of $\sigma=4$ the SSIM of ours and of the reference
model are 0.946 and 0.948, and the high-frequency energy of the non-bone region, relative to the dual-energy soft-tissue image, is 1.09 for
ours and 0.99 for the reference model; LPIPS is 0.234 and 0.166. BS-LDM, which was trained on dual-energy pairs from the same hospital,
reaches 28.57\,dB; it is shown for orientation and is not one of the planned comparisons. Fig.~\ref{fig:pareto} places fidelity against
lesion absorption, and Fig.~\ref{fig:qualitative} shows three held-out cases.

\begin{table*}[!t]
\centering
\caption{Held-out SZCH-X-Rays pairs (51): calibrated fidelity of ours and the comparators (top; HF ratio: high-frequency energy of the
non-bone region relative to the dual-energy soft-tissue image, 1 = equal) and the H2$'$ non-inferiority tests (bottom; paired differences,
image bootstrap, one-sided 95\,\% lower bound). Every method and seed, raw and calibrated, is in Supplementary Table~S1.}
\label{tab:fidelity}
\restable[\textwidth]{tmi_fidelity}
\subtab{H2$'$ non-inferiority}
\restable[\textwidth]{tmi_h2}
\end{table*}

\begin{figure}[!t]
\centering
\safefig[\columnwidth]{figures/pareto.pdf}
\caption{Lesion absorption on VinDr-RibCXR against bone-region PSNR$_{\mathrm{cal}}$ on the 51 held-out SZCH-X-Rays pairs. Line: the
SwinV2-T family at twin weights $w=0\ldots64$ (single seeds light, seed means dark). Markers: ours and the no-twin control
(ConvNeXt-V2-T, $K=3$), the reference model, the v1 model and the identity. Lower absorption at equal fidelity is better.}
\label{fig:pareto}
\end{figure}

\begin{figure*}[!t]
\centering
\safefig[0.8\textwidth]{figures/qualitative_szch.png}
\caption{Held-out SZCH-X-Rays cases (CC BY-NC 4.0 \cite{sun2026bsldm}): radiograph, dual-energy soft-tissue reference, the soft-tissue
estimates of ours and of the reference model, and the bone estimate of ours. The three cases are not hand-picked: they are the held-out
images at the 25th, 50th and 75th percentile of the bone-region PSNR$_{\mathrm{cal}}$ of ours.}
\label{fig:qualitative}
\end{figure*}

\subsubsection{Dose-response of the twin regulariser (H1$'$)}
In the SwinV2-T family (3, 2, 2, 4, 2 and 2 seeds at $w=0,4,8,16,32,64$; Table~\ref{tab:twin}, Fig.~\ref{fig:dose_response}) the mean
absorption on VinDr-RibCXR falls from 0.183 at $w=0$ to 0.043 at $w=64$. The slope of $A$ on $\log_2(1+w)$ is $-0.0233$
$[-0.0249,-0.0217]$; its upper bound is below zero, so H1$'$ passes. On the Qure.ai test split with exact truth the slope is $-0.0174$
$[-0.0190,-0.0160]$, with $A$ levelling off above $w=16$ (0.118, 0.122 and 0.112 at $w=16,32,64$). Bone-region PSNR$_{\mathrm{cal}}$ on the
held-out pairs stays between 21.22 and 21.28\,dB up to $w=16$ and drops to 20.92 and 21.01\,dB at $w=32$ and 64, which is why $w=16$ is the
working point (Fig.~\ref{fig:pareto}).

\begin{table*}[!t]
\centering
\caption{Twin regulariser. Top: dose-response of the lesion absorption on the twin weight (SwinV2-T family): seeds per weight, mean $A$ on
VinDr-RibCXR and on the Qure.ai test split, and the slope with its image-cluster bootstrap interval; each weight averages all of its seeds
(four at $w=16$, seeds 0--3), whereas the encoder comparisons use seeds 0--2. Bottom: ablation at $w=16$ (SwinV2-T, 300k samples, mean over
seeds), one factor changed per row; development evidence (audit on all 75 SZCH images, bone PSNR$_{\mathrm{cal}}$ on the 24 development
pairs).}
\label{tab:twin}
\restable[\textwidth]{dose_response}
\subtab{Ablation at $w=16$}
\restable[\textwidth]{tmi_ablation}
\end{table*}

\begin{figure}[!t]
\centering
\safefig[\columnwidth]{figures/dose_response.pdf}
\caption{Dose-response of the lesion absorption $A$ on VinDr-RibCXR on the twin weight $w$ (SwinV2-T), one marker per seed, with the slope
on $\log_2(1+w)$.}
\label{fig:dose_response}
\end{figure}

\subsubsection{Ablation of the regulariser}
The ablation (Table~\ref{tab:twin}, bottom) is development evidence: it uses the audit on all 75 SZCH images and the 24 development pairs.
At $w=16$ the full regulariser gives $A=0.067$ against 0.190 without it. Three variants lose the effect: a stop-gradient on the branch
without the pasted lesion (0.185), a consistency term over the whole image instead of the lesion support (0.197), and lesion insertion
without the twin term (CLIA, 0.191). Dilating the support by 31\,px keeps part of it (0.137); up to four lesions per twin instead of up to
two keep all of it (0.074). The effect therefore needs gradients through both branches and a consistency term confined to where the lesion
was pasted; showing the model lesions during training is not enough.

\subsubsection{The v1 hypothesis H1}
The v1 hypothesis H1 was tested on synthetic LIDC test CTs with exact truth, three seeds per arm (Supplementary Table~S8), and passes: CLIA
absorbs less than the recipe arm ($-0.022$ $[-0.032,-0.013]$, Holm $p<0.001$) with non-inferior bone MAE. The lesion-weighted loss also
lowers absorption ($-0.019$) and is non-inferior on the test CTs, although it missed the margin on the validation CTs that decided the
selection; combining both lowers absorption most ($-0.066$) but misses the margin. The 300k-sample finals of CLIA and recipe replicate the
contrast ($-0.024$ $[-0.034,-0.014]$). On real radiographs the v1 models remain far above the v2 models ($A=0.289$ and 0.308 on
VinDr-RibCXR; Table~\ref{tab:audit_means}).

% ---------------------------------------------------------------------------------------------------------------------
\subsection{Counterfactual Audit, Rib Benchmark and Real-Nodule Contrast (Study~I: H3\texorpdfstring{$'$}{'},
H4\texorpdfstring{$'$}{'}a)}
\label{sec:res-audit}

\subsubsection{Audit on VinDr-RibCXR (H3$'$)}
On 1,470 sites in 245 VinDr-RibCXR images (Tables~\ref{tab:audit_means} and~\ref{tab:audit}) ours absorbs $A=0.051$ $[0.045,0.057]$ of an
inserted lesion, the reference model 0.149 $[0.138,0.160]$ and the no-twin control 0.178 $[0.167,0.188]$. The paired differences are
$-0.098$ $[-0.108,-0.089]$ against the reference model and $-0.127$ $[-0.134,-0.119]$ against the no-twin control (Holm-adjusted $p<0.001$
for both), so H3$'$ passes. The advantage does not depend on rib overlap here: the interaction between sites with at least 15\,\% overlap
with an annotated rib and the other sites is $0.003$ $[-0.017,0.027]$ against the reference model. On JSRT (secondary) the pattern repeats
($A$ 0.048, 0.118 and 0.121; differences $-0.070$ and $-0.073$) with a larger advantage at bone-overlapping sites (interaction $-0.060$
$[-0.074,-0.046]$); there the overlap was defined by the reference model's bone map. The SZCH audit (development evidence) agrees ($A$ 0.054,
0.130 and 0.195). Fig.~\ref{fig:audit_maps} shows one site.

\begin{table}[!t]
\centering
\caption{Counterfactual lesion audit on VinDr-RibCXR (1,470 sites): mean absorption of each model with image-cluster bootstrap intervals
(lower is better; ensembles of seeds 0--2). Every single seed is in Supplementary Table~S2.}
\label{tab:audit_means}
\restable{tmi_audit_means}
\end{table}

\begin{table*}[!t]
\centering
\caption{H3$'$: paired site-level contrasts on the three audit sets, Holm over the two contrasts, and the interaction with bone overlap
($\geq$15\,\% vs.\ $<$15\,\%; top). SZCH rows use the 51 held-out images (306 sites) and remain development evidence because the audit on
all 75 images was used during development; the JSRT overlap comes from the reference model's bone map. Bottom: rib benchmark on
VinDr-RibCXR (metrics only); every method is in Supplementary Table~S3.}
\label{tab:audit}
\restable[\textwidth]{tmi_audit}
\subtab{Rib benchmark: gradient-energy reduction on rib-boundary pixels and on other lung pixels, selectivity and AUC}
\restable[\textwidth]{tmi_rib}
\end{table*}

\begin{figure*}[!t]
\centering
\safefig[0.95\textwidth]{figures/audit_maps.pdf}
\caption{The audit on one held-out SZCH-X-Rays image (CC BY-NC 4.0): a 7\,mm nodule from a LIDC-IDRI test CT is added to the radiograph;
the change of each bone estimate (red = lesion signal moved into the bone image) and the resulting soft-tissue image. Not hand-picked: eight
placements over the bone-overlapping lung field were drawn with a fixed seed and the one with the median $A(\text{ours})-A(\text{reference})$
is shown.}
\label{fig:audit_maps}
\end{figure*}

\subsubsection{Rib benchmark}
Absorption must be read beside suppression strength, because doing nothing gives $A=0$ (Table~\ref{tab:audit}, bottom). The reference model
removes more edge energy, both on the rib boundaries (48\,\% against 30\,\% for ours) and on the other lung pixels (22\,\% against 8\,\%);
its selectivity is 0.264 against 0.220, and the AUC separating rib from non-rib pixels is the same (0.671 and 0.672). Ours suppresses ribs
less strongly, but its absorption is a third of the reference model's while its rib-edge removal is about 60\,\% of it, and its bone-region
fidelity is non-inferior (H2$'$).

\subsubsection{Encoders (exploratory)}
The twin effect replicates across three encoders trained under the same recipe (Table~\ref{tab:encoders}). On VinDr-RibCXR the twin lowers
$A$ by 0.127 for ConvNeXt-V2-T, 0.114 for SwinV2-T and 0.112 for RAD-DINO, and on JSRT by 0.073, 0.061 and 0.064. ConvNeXt-V2-T was chosen
as ``ours'' on development data before the final evaluation (Section~\ref{sec:m-encoders}); RAD-DINO reaches the same absorption ($A=0.053$
on VinDr-RibCXR).

\begin{table}[!t]
\centering
\caption{Twin regulariser per encoder (exploratory): absorption difference to the matched no-twin control and its interaction with bone
overlap; three seeds per model (seeds 0--2), SZCH on its 51 held-out images. The single seeds of the final models are in
Supplementary Table~S4.}
\label{tab:encoders}
\restable{tmi_s_encoders}
\end{table}

\subsubsection{Real nodules without lesion insertion (H4$'$a)}
Real nodules give a mixed answer (Table~\ref{tab:cnr}). On JSRT (154 nodules) the contrast-to-noise ratio rises by 0.041 with ours and by
0.018 with the reference model; the difference is $+0.023$ $[0.008,0.041]$, and $+0.032$ $[0.014,0.052]$ for the 123 nodules with at least
15\,\% bone overlap. On Node21 (1,474 boxes) the reference model raises contrast more: $+0.103$ against $+0.078$, a difference of $-0.025$
$[-0.033,-0.018]$, and $-0.017$ $[-0.026,-0.009]$ for the 1,109 bone-overlapping boxes. Against the no-twin control ours raises real-nodule
contrast in all four comparisons ($+0.007$ to $+0.027$, all intervals above zero; Holm within data set $p<0.001$).

\begin{table*}[!t]
\centering
\caption{Contrast-to-noise change of real nodules on JSRT and Node21 (soft-tissue image minus original), all lesions and lesions with at
least 15\,\% bone overlap; image-cluster bootstrap intervals.}
\label{tab:cnr}
\restable[0.85\textwidth]{tmi_cnr}
\end{table*}

% ---------------------------------------------------------------------------------------------------------------------
\subsection{Downstream Detection and Classification (Study~I: H4\texorpdfstring{$'$}{'}b)}
\label{sec:res-downstream}
Our soft-tissue image used alone does not help downstream models (Table~\ref{tab:downstream}, Fig.~\ref{fig:downstream_forest}). On
NIH ChestX-ray14 the macro-AUROC of the three-seed ensemble is 0.818 with the original image and 0.817 with ours ($-0.001$
$[-0.004,0.003]$, patient-cluster bootstrap); per class, ours is lower for effusion, infiltration and pleural thickening ($-0.004$,
$-0.007$, $-0.013$) and higher for nodule ($+0.011$; Holm-adjusted DeLong $p<0.05$ for these four). The reference model's soft image alone
gives $+0.004$ $[0.002,0.007]$. In detection the change in CPM with ours alone is $+0.001$ $[-0.023,0.023]$ on Node21, $+0.006$
$[-0.008,0.020]$ on TBX11K and $-0.009$ $[-0.016,-0.002]$ on VinBigData.

Added as extra channels, our suppressed images do help. The original plus our soft image raises NIH macro-AUROC by $+0.006$
$[0.003,0.008]$ (0.818 to 0.824). The three-channel input (original, soft and bone image) raises the CPM on Node21 by $+0.028$
$[0.009,0.047]$ (0.743 to 0.771) and on VinBigData by $+0.008$ $[0.002,0.014]$ (0.726 to 0.734). The gain is established for CPM only:
mAP@50 changes by $+0.030$ $[-0.002,0.062]$ on Node21 and $-0.001$ $[-0.017,0.012]$ on VinBigData. On TBX11K the three-channel input of ours
gives $+0.008$ $[-0.000,0.019]$, whereas the reference model's soft image alone gains $+0.029$ $[0.016,0.043]$. No difference in differences
of the recall at two false positives per image between bone-overlapping and other lesions excludes zero (Node21 $-0.005$ $[-0.071,0.061]$,
VinBigData $+0.009$ $[-0.034,0.053]$ for our soft image); on TBX11K it cannot be estimated, because every test box overlaps bone by at least
20\,\%.

\begin{table*}[!t]
\centering
\caption{Downstream models. Top: classification on NIH ChestX-ray14, macro-AUROC of the three-seed ensemble, per-seed mean $\pm$ SD and
the difference to the original image (patient-cluster bootstrap). Bottom: detection on Node21, TBX11K and VinBigData, CPM and mAP@50 (mean
$\pm$ SD over three seeds), differences to the original image and the difference in differences (DiD) of recall at two false positives per
image between bone-overlapping and other lesions (image bootstrap). ``Original + ours + bone'': the original, soft-tissue and bone images
as three channels.}
\label{tab:downstream}
\restable[\textwidth]{tmi_down_cls}
\subtab{Detection}
\restable[\textwidth]{tmi_down_det}
\end{table*}

\begin{figure}[!t]
\centering
\safefig[\columnwidth]{figures/downstream_forest.pdf}
\caption{Downstream effects of each input against the original image with 95\,\% intervals: macro-AUROC on NIH ChestX-ray14, CPM on the
detection sets.}
\label{fig:downstream_forest}
\end{figure}

% ---------------------------------------------------------------------------------------------------------------------
\subsection{Fusion Study: Overall Performance, Focal Findings and Localisation (Study~II: H6--H8, H10)}
\label{sec:res-fusion}

\subsubsection{Overall performance (H6, H7)}
All 175 runs finished; Table~\ref{tab:fusion} gives the test metrics of the seed ensembles. The hybrid did not improve on the radiograph
alone on any data set (Fig.~\ref{fig:cls_forest}). On NIH the macro-AUROC difference was $+0.0001$ $[-0.0022,0.0024]$ (Holm $p=0.967$), so
a gain larger than 0.0024 is not compatible with the data; on VinBigData it was $+0.0024$ $[-0.0004,0.0053]$ (Holm $p=0.28$) and on
BDCXR-3257 $-0.0042$ $[-0.0108,0.0021]$ (Holm $p=0.39$). The AUPRC differences agreed in direction and none excluded zero. The hybrid did not
differ from the capacity control on any data set (Holm $p=1$). Against channel stacking it was nominally worse on NIH ($-0.0031$
$[-0.0055,-0.0007]$, Holm $p=0.066$; AUPRC $-0.0055$ $[-0.0100,-0.0012]$) and no different elsewhere. Early fusion had the highest NIH
macro-AUROC of the five configurations (0.8273 against 0.8241 for the radiograph alone); this contrast was not one of the tested
hypotheses and is descriptive. Three-stream fusion therefore did not pay for its three times larger encoder computation, and lightweight
channel stacking did at least as well.

\begin{table*}[!t]
\centering
\caption{Fusion study. Top: test macro-AUROC and macro-AUPRC of the seed ensembles (five seeds; BDCXR-3257: out-of-fold over five folds);
Original: radiograph; Soft: soft-tissue image; Early: three images as channels; Hybrid: three streams; Control: radiograph in three
streams. Bottom: H6 (primary) and H7, differences of the seed ensembles with 95\,\% cluster-bootstrap intervals (\num{10000} resamples);
Holm over the three data sets (H6) and over the six tests of H7.}
\label{tab:fusion}
\restable[\textwidth]{tmi_macro}
\subtab{H6 and H7}
\restable[\textwidth]{tmi_primary}
\end{table*}

\begin{figure*}[!t]
\centering
\safefig[0.62\textwidth]{figures/fig_hypothesis_forest_densenet121.pdf}
\caption{Contrasts of the fusion study with 95\,\% cluster-bootstrap intervals. (a) H6 and H7: macro-AUROC differences per data set.
(b) H8: differences in the mean AUROC of each finding group's classes. Filled markers: Holm $p<0.05$; right: Holm-adjusted $p$.}
\label{fig:cls_forest}
\end{figure*}

\subsubsection{Focal and diffuse findings (H8)}
The effect of the inputs depended on the finding group (Table~\ref{tab:groups}, Fig.~\ref{fig:cls_forest}b). On NIH the soft-tissue image
raised the mean AUROC of the focal findings by $+0.0089$ $[0.0040,0.0140]$ (Holm $p=0.0024$) and the hybrid by $+0.0063$ $[0.0005,0.0117]$
(Holm $p=0.101$). Diffuse findings lost with both inputs (hybrid $-0.0052$ $[-0.0077,-0.0027]$, Holm $p=0.0024$; soft $-0.0032$
$[-0.0057,-0.0008]$, Holm $p=0.039$), and boundary findings did not change. On VinBigData the focal estimates had the same direction and
size (hybrid $+0.0085$ $[-0.0023,0.0196]$; soft $+0.0083$ $[-0.0044,0.0202]$) but wide intervals, the diffuse estimates were close to zero
(hybrid $+0.0011$, soft $+0.0003$), and no group difference survived the correction.

\begin{table}[!t]
\centering
\caption{H8: difference in the mean per-class AUROC of each finding group, with 95\,\% cluster-bootstrap intervals; Holm within data set
over groups and contrasts.}
\label{tab:groups}
\restable{tmi_groups}
\end{table}

Per class (Supplementary Table~S9, Fig.~\ref{fig:cls_per_class}), two NIH classes differed between the hybrid and the radiograph after
correction: nodule gained ($+0.0129$, Holm $p=0.00121$) and infiltration lost ($-0.0160$, Holm $p<0.001$; also $-0.0128$ against channel
stacking). Mass did not change ($-0.0004$). No class differed on VinBigData or BDCXR-3257.

\begin{figure*}[!t]
\centering
\safefig[0.7\textwidth]{figures/fig_auroc_per_class_densenet121.pdf}
\caption{NIH ChestX-ray14: test AUROC per class and per finding group for each configuration, and each configuration minus the
radiograph, with 95\,\% patient-cluster bootstrap bands from 1,000 resamples (display only; the bootstrap tests of
Tables~\ref{tab:fusion} and~\ref{tab:groups} use \num{10000}).}
\label{fig:cls_per_class}
\end{figure*}

\subsubsection{Where the streams point (H10)}
Across 984 targets (164 focal) on 880 test images, Grad-CAM++ maps of the radiograph stream hit focal targets less often than other targets
(0.213 against 0.374), whereas soft-tissue stream maps hit both about equally (0.385 against 0.401) (Table~\ref{tab:pointing};
per class in Supplementary Fig.~S7). The soft-minus-radiograph difference was therefore larger for focal targets ($+0.172$) than for the others
($+0.027$): interaction $+0.145$ $[0.092,0.199]$, $p=0.0002$, so H10 passes. Integrated Gradients ranked the two streams the same way
(focal: 0.580 soft-tissue against 0.505 radiograph; other: 0.621 against 0.645; maps in Supplementary Fig.~S6). Bone-stream maps rarely hit
any target (0.049 and 0.167). Fig.~\ref{fig:cls_maps} shows examples.

\begin{table}[!t]
\centering
\caption{H10: pointing-game hit rates on the NIH expert boxes, means over targets of the per-target means over the five hybrid models;
tolerance 15\,px at 256\,px. Last row: soft-tissue minus radiograph stream (Grad-CAM++) and the interaction with its 95\,\%
image-cluster bootstrap interval.}
\label{tab:pointing}
\restable{v3_pointing}
\end{table}

\begin{figure*}[!t]
\centering
\safefig[0.87\textwidth]{figures/fig_xai_comparison_densenet121.pdf}
\caption{Attribution maps of one hybrid model (seed 0) on NIH ChestX-ray14 test images: the three inputs, Grad-CAM++ per stream and the
decoder attention, with the expert boxes. Rows: targets (row label: the box class).}
\label{fig:cls_maps}
\end{figure*}

% ---------------------------------------------------------------------------------------------------------------------
\subsection{External Validation on Paediatric Radiographs (BDCXR-3257: E1--E3)}
\label{sec:res-bdcxr}
On 3,257 paediatric radiographs from a hospital in Dhaka (Table~\ref{tab:bdcxr}), five-fold cross-validated pneumonia classification
(normal, bacterial, viral) reaches a pooled out-of-fold macro-AUROC of 0.814 on the original image and 0.816 on the soft-tissue image of
ours. The difference, $+0.002$ $[-0.005,0.010]$, meets the non-inferiority margin of 0.01 and shows no superiority (E1); the reference
model's soft image gives $+0.001$ $[-0.007,0.010]$ and CLAHE $-0.006$ $[-0.012,0.001]$, which misses the margin. In the lesion audit on 300
of these images (E2) ours absorbs $A=0.084$, the reference model 0.217 and the no-twin control 0.251 (differences $-0.134$
$[-0.144,-0.124]$ and $-0.167$ $[-0.175,-0.160]$), with a larger advantage at bone-overlapping sites. The inserted lesions come from adult
CTs, so this is a stress test rather than a clinical finding. The ensemble disagreement on these children (E3) is 2.7 times that on the
adult SZCH-X-Rays pairs (median lung-field spread 0.0076 against 0.0028) and is larger in younger children (Spearman with age $-0.22$).

\begin{table*}[!t]
\centering
\caption{BDCXR-3257. Top (E1): pneumonia classification (normal, bacterial, viral), five folds over near-duplicate groups $\times$ three
seeds; pooled out-of-fold macro-AUROC, per-fold mean $\pm$ SD, the difference to the original image (group bootstrap), per-class AUROC and
balanced accuracy. Bottom (E2): lesion audit on a fixed class-stratified subset of 300 images (1,800 sites).}
\label{tab:bdcxr}
\restable[\textwidth]{tmi_s_bdcxr_cls}
\subtab{Lesion audit (E2)}
\restable[0.4\textwidth]{tmi_s_bdcxr_audit}
\end{table*}

% ---------------------------------------------------------------------------------------------------------------------
\subsection{Uncertainty of the Suppressor and the Classifiers (H5\texorpdfstring{$'$}{'}, H9)}
\label{sec:res-uq}

\subsubsection{Reliability of the bone-image uncertainty (H5$'$)}
On the Qure.ai test split with exact truth the split-conformal intervals of the three-model ensemble cover 0.900 $[0.897,0.903]$ of the
pixels at a nominal 0.90, and the ensemble spread ranks the absolute bone error well (Spearman 0.76; Table~\ref{tab:reliability};
per-image coverage in Supplementary Fig.~S1). On the 51 held-out SZCH-X-Rays pairs, measured against a pseudo-truth bone image from the dual-energy pair,
the same intervals cover only 0.139 $[0.134,0.145]$ and the spread ranks the error weakly (Spearman 0.11 $[0.07,0.15]$). The uncertainty is
calibrated in the training domain and does not transfer to real radiographs without recalibration.

\begin{table*}[!t]
\centering
\caption{Bone-image uncertainty. Top (H5$'$): ensemble spread and split-conformal intervals ($K=3$, nominal coverage 0.90), with
image-bootstrap intervals. Bottom (H9, exploratory and post hoc): intervals recalibrated on the 24 SZCH-X-Rays development pairs and
tested on the 51 held-out pairs at nominal 0.90; Mondrian: calibration per stratum of predicted bone overlap and ensemble spread.}
\label{tab:reliability}
\restable[\textwidth]{reliability}
\subtab{Recalibration on real pairs (H9)}
\restable[\textwidth]{v3_bone_intervals}
\end{table*}

\subsubsection{Recalibration on real pairs (H9, post hoc)}
An exploratory analysis, added after this result was known and therefore post hoc on already-seen pairs, recalibrated the intervals on the
24 SZCH-X-Rays development pairs and tested them on the 51 held-out pairs (Table~\ref{tab:reliability}, bottom; Supplementary Fig.~S2).
Pooled coverage was 0.886 with one global calibration and 0.895 with a separate calibration per stratum of predicted bone overlap and
ensemble spread (Mondrian), whose intervals were less than half as wide (mean width 0.215 against 0.447). A few real pairs thus restore
nominal coverage on average, while the coverage of single images still varies.

\subsubsection{Classification uncertainty}
Conformal risk control held the NIH false-negative rate at 0.060--0.069 on the test split for nominal 0.10, but the label sets held
6.33--6.78 of 14 labels on average (Supplementary Table~S10 and Fig.~S4). On BDCXR-3257 the Mondrian sets covered 0.897--0.900 of
the images at nominal 0.90 with 1.76--1.81 of 3 classes per set. In selective classification (Supplementary Fig.~S5), adding the
suppressor's spread to the confidence score raised the AURC in every configuration: NIH 0.0352--0.0358 with the entropy alone against
0.0372--0.0375 with the spread, both on the \num{25580} test images with a lung-field spread (16 of \num{25596} have an empty lung mask),
and BDCXR-3257 0.1754--0.1938 against 0.2138--0.2260. Every single run is shown next to the seed ensembles in Supplementary Fig.~S3.

% ---------------------------------------------------------------------------------------------------------------------
\subsection{Site-Appearance Calibration (Exploratory, H16--H18)}
\label{sec:res-sitecal}
\begin{table*}[!t]
\centering
\caption{Site-appearance calibration (exploratory). Top: fidelity on the 51 held-out SZCH-X-Rays pairs, every method as is and with the
smooth correction of Eq.~\eqref{eq:sitecal} fitted on the 24 development pairs (BS-LDM, trained on this site's pairs, is not corrected).
Middle: H16--H18, paired differences with image bootstrap (H16, H16b, H17: one-sided 95\,\% lower bounds) or image-cluster bootstrap
intervals (absorption on 1,470 VinDr-RibCXR sites), \num{10000} resamples. Bottom: VinDr-RibCXR (metrics only), mean lesion absorption of
one audit run with all four models (shared sites) and the rib benchmark, as is and site-calibrated.}
\label{tab:sitecal}
\restable[\textwidth]{site_cal_fidelity}
\subtab{H16--H18}
\restable[\textwidth]{site_cal_tests}
\subtab{VinDr-RibCXR: absorption and rib benchmark}
\restable[\textwidth]{site_cal_ribcxr}
\end{table*}

The smooth site correction of Section~\ref{sec:m-sitecal}, fitted on the 24 development pairs, raised the bone-region PSNR$_{\mathrm{cal}}$
of ours on the 51 held-out pairs from 21.45 to 23.58\,dB, a paired gain of 2.12\,dB (one-sided 95\,\% lower bound 1.99\,dB; H16 passes),
and its whole-image PSNR$_{\mathrm{cal}}$ from 25.56 to 28.30\,dB (gain 2.74\,dB, lower bound 2.61\,dB; H16b passes;
Table~\ref{tab:sitecal}). Every CT-supervised model gained about as much (reference model 21.51 to 23.92\,dB, v1 model 21.95 to 24.06\,dB,
no-twin control 21.45 to 23.74\,dB) and the unprocessed radiograph less (20.21 to 21.46\,dB): the correction calibrates the site rather
than improving a suppressor, and the order of the CT-supervised models did not change. After calibration ours remained non-inferior to the
calibrated reference model, by a narrower margin than before ($-0.35$\,dB, lower bound $-0.44$\,dB against the margin of $-0.5$\,dB; H17
passes), and the no-twin control was slightly ahead of ours (23.74 against 23.58\,dB; not tested).

The calibration did not cost lesion fidelity (Table~\ref{tab:sitecal}, bottom): on VinDr-RibCXR the calibrated ours absorbed 0.036 of an
inserted lesion and the calibrated reference model 0.132 (difference $-0.096$ $[-0.105,-0.087]$; H18 passes), and the calibration lowered
the absorption of both models slightly ($-0.015$ and $-0.017$); the uncalibrated models of this run reproduced the audit of
Table~\ref{tab:audit_means}. It weakened suppression a little and made it more selective: the rib-edge reduction of ours fell from 0.298 to
0.282 and the reduction on other lung pixels from 0.078 to 0.050, so the selectivity rose from 0.220 to 0.232. BS-LDM, trained on this
site's dual-energy pairs, remains about 5\,dB ahead in the bone region (28.57 against 23.58\,dB) and 3.65\,dB in the whole image.

